% BHU PYQ -- Transcribed & Corrected
% Paper: ARBMD-21 : Basic Arabic - II
% Exam: Under Graduate Semester II Examination, 2024-25 (NEP)
% Department: Department of Arabic

\documentclass[11pt,a4paper]{article}
\usepackage[a4paper,top=2.5cm,bottom=2.5cm,left=2.8cm,right=2.8cm,headheight=14pt]{geometry}
\usepackage{amsmath,amssymb}
\usepackage[T1]{fontenc}
\usepackage[utf8]{inputenc}
\usepackage{lmodern,microtype,fancyhdr,enumitem,hyperref,lastpage,parskip}

\hypersetup{
    colorlinks=true,
    urlcolor=black,
    linkcolor=black,
    pdftitle={ARBMD21: Basic Arabic II -- BHU UG Sem II 2024-25 (NEP)}
}

\pagestyle{fancy}
\fancyhf{}
\renewcommand{\headrulewidth}{0.4pt}
\renewcommand{\footrulewidth}{0.4pt}
\fancyhead[L]{\small\textit{Banaras Hindu University}}
\fancyhead[R]{\small\textit{ARBMD-21: Basic Arabic - II}}
\fancyfoot[C]{\small Page \thepage\ of \pageref{LastPage}}

\newlist{parts}{enumerate}{2}
\setlist[parts,1]{label=(\roman*),leftmargin=*,itemsep=4pt,topsep=2pt}
\newcommand{\pts}[1]{\hfill\textit{[#1]}}

\begin{document}

\begin{center}
  {\large\bfseries BANARAS HINDU UNIVERSITY}\\[2pt]
  {\normalsize Under Graduate Semester II Examination, 2024-25 (NEP)}\\[6pt]
  \rule{0.8\linewidth}{0.5pt}\\[6pt]
  {\large\bfseries DEPARTMENT OF ARABIC}\\[4pt]
  {\large\bfseries Paper Code: ARBMD-21 : Basic Arabic - II}\\[4pt]
  {\normalsize Time: 3 Hours\hfill Maximum Marks: 60}
\end{center}

\rule{\linewidth}{0.4pt}

\smallskip
\noindent\textit{Instructions: Attempt any FIVE questions. All questions carry equal marks (12 marks each).}

\bigskip

\noindent\textbf{Question 1.} Define and explain the Nominal Sentence (الجملة الاسمية) in Arabic grammar. Describe the role of Subject (المبتدأ) and Predicate (الخبر) with suitable examples.\pts{12}

\medskip\rule{\linewidth}{0.2pt}\medskip

\noindent\textbf{Question 2.} What is a Noun-Adjective Phrase (المركب التوصيفي)? Discuss the four conditions of agreement between the Noun (الموصوف) and its Adjective (الصفة) with examples.\pts{12}

\medskip\rule{\linewidth}{0.2pt}\medskip

\noindent\textbf{Question 3.} Translate the following Arabic sentences into English or Hindi:
\begin{parts}
  \item الكتاب الجديد على الطاولة الكبيرة.
  \item الولد النشيط يدرس في المدرسة.
  \item البيت الجميل قريب من المسجد.
  \item المعلمة الماهرة تشرح الدرس بوضوح.
\end{parts}\pts{12}

\medskip\rule{\linewidth}{0.2pt}\medskip

\noindent\textbf{Question 4.} Explain the Possessive Construction (المركب الإضافي) in Arabic. Detail the rules governing the Possessed (المضاف) and Possessor (المضاف إليه) with illustrative examples.\pts{12}

\medskip\rule{\linewidth}{0.2pt}\medskip

\noindent\textbf{Question 5.} Write short grammatical notes on any TWO of the following:
\begin{parts}
  \item Prepositions (حروف الجر) and their impact on nouns.
  \item Demonstrative Pronouns (أسماء الإشارة) for near and far objects.
  \item Attached Personal Pronouns (الضمائر المتصلة).
\end{parts}\pts{12}

\vfill
\rule{\linewidth}{0.4pt}
\begin{center}\small
\href{https://bhu-exam.vercel.app/}{Click here to visit our website}\\[4pt]
\href{https://me-aryan.vercel.app/}{Click here to visit our contributor}
\end{center}

\end{document}
